%! AP Statistics — Unit 3, Lesson 3.3 — Activity KEY
\documentclass[10pt]{article}
\usepackage{apstats-article}
\usepackage{apstats-key}

\begin{document}

\pageheader{Unit 3, Lesson 3.3}{Group Activity KEY: Design Your Sample}
\namepartnerperiod

\vspace{0.1in}

\begin{scenariobox}[Context]{sky}
\small
A large urban school district has 8{,}000 students across 4 grade levels (9--12),
16 schools, and classes averaging 30 students each. The district wants to survey
students about their satisfaction with the school lunch program.
Your group will design a sampling plan using each method and evaluate which is best.
\end{scenariobox}

\vspace{0.1in}

\textbf{Part 1: Design a sampling plan using each method.}

\begin{practicebox}
\small
\textbf{A. Simple Random Sample (SRS) --- Target: 200 students}

\begin{enumerate}[leftmargin=1.5em, itemsep=6pt]
  \item How would you number the students? \makebox[9cm][l]{\ans{Give each a unique number 1--8{,}000.}}
  \item What device would you use to select 200? \makebox[9cm][l]{\ans{Calculator randInt or a random number table.}}
  \item One advantage of SRS here: \makebox[9cm][l]{\ans{Every group of 200 has the same chance.}}
  \item One disadvantage: \makebox[9cm][l]{\ans{Could under-represent a grade level by chance.}}
\end{enumerate}
\end{practicebox}

\vspace{0.1in}

\begin{practicebox}
\small
\textbf{B. Stratified Random Sample --- Target: 200 students}

\begin{enumerate}[leftmargin=1.5em, itemsep=6pt]
  \item What are the natural strata in this district? \makebox[8cm][l]{\ans{The four grade levels (9, 10, 11, 12).}}
  \item How many students would you sample from each stratum (be proportional)?\\
        Grade 9: \makebox[2cm][l]{\ans{50}} \quad Grade 10: \makebox[2cm][l]{\ans{50}} \quad Grade 11: \makebox[2cm][l]{\ans{50}}
        \quad Grade 12: \makebox[2cm][l]{\ans{50}}
  \item Why is stratifying by grade level useful here?\\
        \ansline{Each grade is fairly represented; lunch satisfaction may}\\[1pt]\ansline{differ by grade, so stratifying reduces variability.}
\end{enumerate}
\end{practicebox}

\vspace{0.1in}

\begin{practicebox}
\small
\textbf{C. Cluster Sample --- Target: all students in selected schools}

\begin{enumerate}[leftmargin=1.5em, itemsep=6pt]
  \item The district has 16 schools. If you randomly select 4 schools, how many
        students would be in your sample (approximately)? \makebox[4cm][l]{\ans{$4 \times 500 = 2{,}000$}}
  \item Is a school a reasonable cluster? Explain.\\
        \ansline{Yes --- each school has students from all grade levels,}\\[1pt]\ansline{so each cluster looks like a mini-district.}
  \item One practical advantage of cluster sampling here: \makebox[9cm][l]{\ans{Survey whole schools, not scattered students.}}
\end{enumerate}
\end{practicebox}

\vspace{0.1in}

\begin{practicebox}
\small
\textbf{D. Systematic Random Sample --- Target: approximately 200 students}

\begin{enumerate}[leftmargin=1.5em, itemsep=6pt]
  \item The district has 8{,}000 students on a master list. To get 200, your sampling
        interval $k$ is: $k = $ \makebox[2cm][l]{\ans{40}}
  \item How do you choose your random start? \makebox[9cm][l]{\ans{Random number 1--40, then every 40th.}}
  \item A potential problem with systematic sampling here: \makebox[9cm][l]{\ans{A patterned list could over-represent some grades.}}
\end{enumerate}
\end{practicebox}

\vspace{0.1in}

\textbf{Part 2: Compare and recommend.}

\begin{practicebox}
\small
\begin{enumerate}[leftmargin=1.5em, itemsep=8pt]
  \item Which method best ensures that each grade level is fairly represented?
        \makebox[4cm][l]{\ans{Stratified}} \quad Explain: \makebox[7cm][l]{\ans{Samples every grade proportionally.}}
  \item Which method would be easiest to implement in a school setting?
        \makebox[4cm][l]{\ans{Cluster}} \quad Explain: \makebox[7cm][l]{\ans{Survey a few whole schools.}}
  \item Which method would your group recommend for this district? Justify in 2--3
        sentences.\\
        \ansline{Stratified by grade level: it guarantees each grade is}\\[1pt]\ansline{proportionally represented and allows comparing grades.}\\[1pt]\ansline{Cluster sampling is practical but may miss grade-level differences.}
\end{enumerate}
\end{practicebox}

\end{document}
